\section{Task Definition}

\paragraph{Input.}
Each example contains a paper title and optional abstract, one reviewer concern, author response candidates, an aligned response excerpt, and a revision summary.
The unit of prediction is a single issue, not a paper-level decision.
This design prevents a model from hiding behind a global paper-quality judgment when one concern was fixed and another remains unresolved.

\paragraph{Output labels.}
The label set is:
\textit{fixed}, \textit{partially fixed}, \textit{unresolved}, and \textit{regressed}.
\textit{Fixed} means the exact concern is addressed with concrete revision evidence.
\textit{Partially fixed} means the revision makes real progress but leaves a material part open.
\textit{Unresolved} means the concern remains present, is deferred, or is reframed without enough evidence.
\textit{Regressed} means the attempted fix introduces a new issue on the same axis or makes the paper less reliable.

\paragraph{Adjudication protocol.}
To reduce ambiguity at the fixed/partial/unresolved boundary, we apply the same decision order for every row.
\textbf{(1) Axis match:} check whether the revision evidence addresses the same technical axis as the reviewer concern (not just tone alignment).
\textbf{(2) Materiality:} check whether the highest-impact part of the concern is resolved for scientific judgment (for example, a missing core comparison vs a minor wording tweak).
\textbf{(3) Evidence sufficiency:} require concrete revision-facing evidence (new experiment, corrected claim, revised method detail, released artifact, or explicit limitation narrowing) rather than acknowledgement-only responses.
This yields \textit{fixed} when all three checks pass, \textit{partially fixed} when axis match and partial material progress exist but a material gap remains, \textit{unresolved} when evidence is insufficient or deferred, and \textit{regressed} when the attempted fix degrades reliability on the same axis.

\paragraph{Boundary examples in main narrative.}
Two recurring boundary patterns illustrate this protocol.
\textbf{Fixed vs partially fixed:} adding one ablation requested by the reviewer can be \textit{partially fixed} if the core missing comparison remains absent, and only becomes \textit{fixed} once the requested core comparison is present with revision evidence.
\textbf{Partially fixed vs unresolved:} a detailed rebuttal that promises future experiments remains \textit{unresolved} without revision evidence, while a partial empirical addition that narrows but does not close the concern is \textit{partially fixed}.
Appendix Table~\ref{tab:worked-issue-examples} provides compact concern/evidence snippets for all four labels.

\paragraph{Evaluation.}
Macro-F1 is the primary metric because label skew is central to the task.
Accuracy is reported only with null baselines and per-label recovery.
We explicitly track fixed-case and unresolved-case F1 because stale criticism and over-crediting are the most consequential failures for review support.
